\chapter{Construcción de hardware y comunicación con ROS 2}
\label{chap:robot_real}

Este capítulo describe la arquitectura físico-electrónica del robot diferencial y su integración con \gls{ROS2} Foxy sobre una Raspberry Pi 4. Se presentan los dominios de potencia y comunicación, el esquema general de conexiones, la integración con \texttt{ros2\_control}, el controlador diferencial y el flujo de datos entre la capa de alto nivel y el firmware del \gls{Arduino}. Los detalles operativos de cableado, reglas \texttt{udev}, QoS y configuración extendida se trasladan al \ref{anx:hardware_ros2}.

%------------------------------------------------------------------
\section{Arquitectura física y de datos}

La plataforma se organiza en dos dominios principales. El dominio de potencia utiliza un bus de 12~V para alimentar los motores DC mediante el driver L298N, y un riel regulado de 5~V para alimentar la Raspberry Pi, el \gls{Arduino}, el \gls{LIDAR} y los periféricos. Esta separación reduce la interferencia eléctrica producida por los motores sobre la electrónica de control.

El dominio de comunicaciones conecta la Raspberry Pi con el \gls{Arduino} mediante enlace serial USB, utilizado para enviar referencias de velocidad y recibir datos de encoders. El sensor RPLidar A1 entrega barridos láser publicados en \texttt{/scan}, empleados por \texttt{slam\_toolbox} y por los mapas de costo de Nav2. La Figura~\ref{fig:esquema_conexiones} resume las conexiones eléctricas y de datos del prototipo.

\utemFigurePropia{2.Texto/Imag/circuit_diagram_bb.png}{0.90\textwidth}
{Esquema de conexiones eléctricas y de datos del robot diferencial.}
{fig:esquema_conexiones}

%------------------------------------------------------------------
\section{Integración con ROS 2}

En hardware real, el archivo \texttt{launch\_robot.launch.py} inicia los nodos necesarios para publicar la descripción del robot, activar \texttt{ros2\_control} y cargar el controlador diferencial. La arquitectura incluye:

\begin{itemize}
\item \texttt{robot\_state\_publisher}, encargado de publicar las transformaciones del robot.
\item \texttt{controller\_manager}, que gestiona la interfaz de hardware y los controladores.
\item \texttt{diff\_drive\_controller}, que convierte comandos de velocidad en referencias para las ruedas.
\item \texttt{joint\_state\_broadcaster}, que publica el estado de las articulaciones.
\end{itemize}

La cadena de transformaciones mantiene la misma convención usada en simulación:

\begin{center}
\texttt{map} $\rightarrow$ \texttt{odom} $\rightarrow$ \texttt{base\_link} $\rightarrow$ \texttt{laser\_frame}
\end{center}

En esta cadena, \texttt{slam\_toolbox} publica la transformación \texttt{map}$\rightarrow$\texttt{odom}, mientras que \texttt{diff\_drive\_controller} publica \texttt{odom}$\rightarrow$\texttt{base\_link} a partir de la odometría de encoders.

%------------------------------------------------------------------
\section{Teleoperación y multiplexación de comandos}

El nodo \texttt{twist\_mux} arbitra las fuentes de velocidad antes de enviar comandos al controlador diferencial. La teleoperación tiene mayor prioridad que la navegación autónoma y funciona como mecanismo de seguridad mediante un \textit{dead-man switch}. Su salida se remapea hacia \texttt{/diff\_cont/cmd\_vel\_unstamped}, tópico de entrada del controlador diferencial.

\begin{table}[H]
\centering
\caption{Fuentes de velocidad gestionadas por \texttt{twist\_mux}.}
\label{tab:twist_mux_resumen}
\small
\begin{tabular}{lcc}
\toprule
\textbf{Fuente} & \textbf{Tópico} & \textbf{Prioridad} \\
\midrule
Navegación autónoma & \texttt{/cmd\_vel} & Baja \\
Teleoperación & \texttt{/cmd\_vel\_joy} & Alta \\
\bottomrule
\end{tabular}
\end{table}

%------------------------------------------------------------------
\section{Control diferencial e interfaz de hardware}
\label{sec:diff_drive_params}

El controlador diferencial recibe comandos de velocidad lineal y angular, y los convierte en velocidades de rueda mediante el modelo cinemático descrito en el Capítulo~\ref{marco_teorico}. Los parámetros geométricos son consistentes con el modelo físico y el gemelo digital: \texttt{wheel\_radius} = 0.0335~m y \texttt{wheel\_separation} = 0.188~m.

\begin{table}[H]
\centering
\caption{Parámetros principales del controlador diferencial.}
\label{tab:diff_drive_params}
\small
\begin{tabular}{lll}
\toprule
\textbf{Parámetro} & \textbf{Valor} & \textbf{Descripción} \\
\midrule
\texttt{linear.x.max\_velocity} & $\pm$0.13~m/s & Velocidad lineal máxima \\
\texttt{angular.z.max\_velocity} & $\pm$0.4~rad/s & Velocidad angular máxima \\
\texttt{update\_rate} & 30~Hz & Frecuencia del controlador \\
\texttt{enable\_odom\_tf} & true & Publica \texttt{odom}$\rightarrow$\texttt{base\_link} \\
\bottomrule
\end{tabular}
\end{table}

La comunicación con el firmware se realiza mediante el plugin \texttt{DiffDriveArduino}, que traduce las velocidades angulares de rueda hacia el formato de \textit{ticks/frame} utilizado por el \gls{Arduino}. A su vez, recibe los conteos de encoder y los convierte en posición y velocidad de cada rueda para alimentar la odometría.

%------------------------------------------------------------------
\section{Flujo de datos del sistema}
\label{sec:flujo_datos}

El flujo de control entre \gls{ROS2}, el plugin de hardware y el firmware se resume en las siguientes etapas:

\begin{enumerate}
\item Un nodo de navegación o teleoperación publica un comando de velocidad.
\item \texttt{twist\_mux} selecciona la fuente de mayor prioridad.
\item \texttt{diff\_drive\_controller} convierte el comando $(v,\omega)$ en velocidades de rueda.
\item \texttt{DiffDriveArduino} traduce esas velocidades a \textit{ticks/frame}.
\item El \gls{Arduino} actualiza las referencias PID y aplica \gls{PWM} a los motores.
\item El plugin solicita periódicamente los encoders y actualiza la odometría.
\item \texttt{diff\_drive\_controller} publica \texttt{/odom} y la transformación correspondiente.
\end{enumerate}

El firmware incorpora una auto-detención de seguridad: si no recibe comandos durante más de 2~s, detiene los motores. Esta protección reduce el riesgo ante fallas de comunicación, reinicio de nodos o pérdida de conexión.

%------------------------------------------------------------------
\section*{Conclusión del capítulo}

La arquitectura de hardware integra una capa de control embebido en \gls{Arduino} con una capa de alto nivel en \gls{ROS2}, ejecutada sobre Raspberry Pi 4. La comunicación mediante \texttt{ros2\_control}, \texttt{diff\_drive\_controller} y el plugin \texttt{DiffDriveArduino} permite cerrar el ciclo entre comandos de navegación, control PID de ruedas y publicación de odometría. Esta base permite implementar los módulos de SLAM, navegación autónoma y exploración descritos en los capítulos siguientes.